\Needspace{13\baselineskip}
\chapterrule
\section{训练得很好，为什么新题反而做错？}
\label{intro:18-generalization}
\chaptercompass{训练集很准，为何新样本仍会出错？}{用独立验证集区分拟合训练样本与泛化能力。}{训练/验证/测试样本 $\to$ 各自损失}{模型参数在训练集更新；验证集不反传}
\modelpipelinefocus{0}{泛化与验证}



\subsection{欠拟合、合适的拟合、过拟合}

真实规律接近 $y=2x+1$，十个测量样本带有噪声。
直线能抓住总体趋势；高次多项式穿过所有点，也会把噪声学进去。
\textbf{欠拟合}（underfitting）是训练和新数据都表现不好；
\textbf{合适的拟合}抓住主要模式；
\textbf{过拟合}（overfitting）则是训练表现优异、新数据表现差。
\textbf{泛化能力}就是模型面对训练时没见过的数据的表现。

% Same five observations and exact functions as the accompanying hand calculation.
\noindent\begin{minipage}{\linewidth}
\tagstructbegin{tag=Figure,alttext={三个面板使用相同数据：实心训练点为负一与负一点五、零与二、一与二点五；空心验证点为负二分之一与八分之一、二分之一与十七分之八。依次画常数一、直线二倍 x 加一、二次曲线负一点五倍 x 平方加二倍 x 加二。二次曲线穿过训练点，却在两个区间内验证点上比直线误差更大。}}
\begin{center}
\begin{tikzpicture}[x=15mm,y=8mm]
\foreach \dx/\lab/\formula in {0/欠拟合/{$\widehat{y}=1$},3.1/直线拟合/{$\widehat{y}=2x+1$},6.2/曲线拟合/{$\widehat{y}=-1.5x^2+2x+2$}}{
  \begin{scope}[shift={(\dx,0)}]
    \draw[->,linegray] (-1.18,0)--(1.25,0) node[right,font=\scriptsize] {$x$};
    \draw[->,linegray] (0,-1.9)--(0,3.4) node[above,font=\scriptsize] {$y$};
    \foreach \px in {-1,1}{
      \draw[linegray] (\px,.08)--(\px,-.08);
      \node[below,font=\scriptsize] at (\px,-.08) {$\px$};
    }
    \node[font=\small\bfseries] at (0,4.3) {\lab};
    \node[font=\scriptsize] at (0,3.7) {\formula};
  \end{scope}
}
\draw[flowcolor,very thick] (-1.05,1)--(1.05,1);
\begin{scope}[shift={(3.1,0)}]
  \draw[flowcolor,very thick] plot[domain=-1.05:1.05,samples=2] (\x,{2*\x+1});
\end{scope}
\begin{scope}[shift={(6.2,0)}]
  \draw[flowcolor,very thick] plot[domain=-1.05:1.05,samples=40] (\x,{-1.5*\x*\x+2*\x+2});
\end{scope}
\foreach \dx in {0,3.1,6.2}{
  \begin{scope}[shift={(\dx,0)}]
    \foreach \px/\py in {-1/-1.5,0/2,1/2.5}
      \fill[noisecolor] (\px,\py) circle(2.3pt);
    \foreach \px/\py in {-.5/.125,.5/2.125}
      \draw[noisecolor,thick,fill=white] (\px,\py) circle(2.6pt);
  \end{scope}
}
\end{tikzpicture}

{\footnotesize 黑点是带测量噪声的训练样本，蓝灰线是模型预测；空心点是两个验证样本，数据与正文相同。}
\captionof{figure}{用相同三点训练、另两点验证：穿过训练点的曲线，未必更适合新题}\label{fig:fitting}
\end{center}
\tagstructend
\end{minipage}\par


\subsection{为什么需要训练、验证和测试三份数据？}

\begin{center}
\small
\tagpdfsetup{table/header-rows={1}}
\begin{tabularx}{\linewidth}{@{}p{21mm}X@{}}\toprule
集合 & 用途\\\midrule
训练集 & 用于计算梯度并更新模型参数\\
验证集 & 用于选择超参数、比较模型与决定何时停止\\
测试集 & 留到最后使用，估计模型对新数据的泛化表现\\
\bottomrule\end{tabularx}
\end{center}
如果反复根据测试集成绩修改模型，测试集也会被间接“学会”，无法再充当可信的独立评估。数据切分还应避免重复样本和时间泄漏。

\subsection{四种常见的抗过拟合办法}

\textbf{一是增加数据并合理增强。}
适度旋转、平移、裁剪图片可增加训练模式，但增强须保持标签；例如把数字 $6$ 旋转 $180^\circ$，会把它变成另一类的外形。

\textbf{二是控制模型复杂度。}
先用简单模型作基线，再根据验证表现调整容量；泛化要同时考察数据、容量与训练过程。

\textbf{三是正则化。}
把\textbf{参数大小的惩罚}加到原来的任务损失：
\begin{equation}
\begin{aligned}
L_{\rm total}(w_1,\ldots,w_n)
&=L_{\rm data}(w_1,\ldots,w_n)
+\lambda\sum_i|w_i| &&\text{（L1）},\\
L_{\rm total}(w_1,\ldots,w_n)
&=L_{\rm data}(w_1,\ldots,w_n)
+\lambda\sum_iw_i^2 &&\text{（平方 L2 惩罚）}.
\end{aligned}
\end{equation}
$\lambda\ge0$ 控制惩罚力度，称为\textbf{正则化系数}。数值越大，模型越倾向较小的参数，过强时会压低拟合能力。

\textbf{四是提前停止（early stopping）。}
训练损失下降而验证损失变差时，可以保留验证表现最好的检查点。

% The Figure structure intentionally includes the non-floating caption.
\tagstructbegin{tag=Figure,alttext={训练与验证损失曲线不同步，提前停止选择验证误差较低时的参数}}
\begin{center}
\begin{tikzpicture}
\begin{axis}[
 width=112mm,height=46mm,
 xmin=.8,xmax=8.3,ymin=.2,ymax=1.28,
 axis lines=left,xlabel={训练轮次（Epoch）},ylabel={损失},
 xtick={1,2,3,4,5,6,7,8},
 tick label style={font=\footnotesize},
 label style={font=\small},
 legend style={at={(.98,.97)},anchor=north east,font=\footnotesize,
 fill=white,draw=linegray},
]
\addplot[blue,very thick,mark=*] coordinates {
(1,1.17) (2,.89) (3,.71) (4,.55) (5,.44) (6,.35) (7,.29) (8,.25)};
\addlegendentry{训练损失}
\addplot[orange,very thick,mark=square*] coordinates {
(1,1.22) (2,.96) (3,.80) (4,.74) (5,.78) (6,.84) (7,.94) (8,1.06)};
\addlegendentry{验证损失}
\draw[dashed,teal,thick] (axis cs:4,.2)--(axis cs:4,1.2);
\node[teal,anchor=south west,font=\footnotesize] at (axis cs:4.12,.23) {最佳验证轮次};
\end{axis}
\end{tikzpicture}

{\footnotesize 人工示意数据：训练损失持续下降，
验证损失在第 4 轮后反而上升；选检查点应看验证集。}
\end{center}
\begin{center}\captionof{figure}{训练损失与验证损失走势及最优验证轮次}\label{fig:early-stopping}\end{center}
\tagstructend


\subsection{Dropout 到底丢弃了什么？}
\phantomsection\label{concept:dropout}

Dropout 在训练时随机将\textbf{部分中间激活值}暂时置零。
例如 $(2,1,3,4)$ 经两次掩码可分别变为 $(2,0,3,0)$、$(0,1,3,4)$；
反向缩放 Dropout 再将保留值除以保留概率，以控制期望大小。
它屏蔽当前激活，参数仍保留。
推理前用 \texttt{model.eval()} 切到\textbf{评估模式}，关闭 Dropout；
仅关闭梯度记录不会切换模式。

\begin{conclusion}
训练误差与验证误差回答不同问题。独立测试集、正则化和提前停止帮助我们判断学到的模式能否迁移到新数据。
下一部分面对更现实的输入：如果特征不止一个，怎样一起计算？
\end{conclusion}

\chapterchecktasks{18-generalization}{正文基础三题}{
\item \textbf{预测：}训练损失持续下降，验证损失却开始上升，模型更可能在欠拟合、过拟合、还是健康泛化？先提出一种排查办法。\chapterinlinehint{先比较训练曲线与验证曲线的方向。}
\item \textbf{手算：}三轮验证损失为 $(0.6,0.4,0.5)$。按最低验证损失早停，选哪一轮？最低是多少？\chapterinlinehint{在验证记录中找最小值及其所在轮次。}
\item \textbf{找错：}同学说“Dropout 训练时随机删除一些权重参数，再由测试集挑最佳 Dropout 率”。指出两处错误。\chapterinlinehint{分别核对 Dropout 的作用对象与调参数据集。}
}
\chapterrecallanswer{训练误差变小不保证验证误差变小；测试集必须留给最后的独立评估。}




